\documentclass[11pt,a4paper]{article}
\usepackage[utf8]{inputenc}
\usepackage[spanish]{babel}
\usepackage{amsmath, amssymb}
\usepackage[margin=2.5cm]{geometry}
\usepackage{fancyhdr}
\usepackage{tabularx}
\usepackage{booktabs}
\usepackage{tcolorbox}
\usepackage{enumitem}

% Configuración de colores corporativos UCB
\definecolor{UCBblue}{RGB}{0, 51, 153}
\definecolor{UCBgray}{RGB}{100, 100, 100}
\definecolor{UCBred}{RGB}{153, 0, 0}

% Configuración de encabezado y pie de página
\pagestyle{fancy}
\fancyhf{}
\fancyhead[L]{\textcolor{UCBblue}{\textbf{Ingeniería Mecatrónica - UCB Tarija}}}
\fancyhead[R]{\textcolor{UCBgray}{Guión de Clase Ajustado - Sesión 3}}
\fancyfoot[C]{\thepage}
\renewcommand{\headrulewidth}{0.4pt}
\renewcommand{\headrule}{\hbox to\headwidth{\color{UCBblue}\leaders\hrule height \headrulewidth\hfill}}

\begin{document}

\begin{center}
    {\Large \textbf{\textcolor{UCBblue}{PLANIFICACIÓN DIDÁCTICA Y GUIÓN DE CLASE}}}\\
    \vspace{0.3cm}
    {\large \textbf{Sesión 3 (Ajustada): Análisis Fasorial y Dinámica Frecuencial}}\\
    \vspace{0.5cm}
    \textbf{Asignatura:} Taller de Nivelación / Introducción a la Mecatrónica \\
    \textbf{Docente:} M.Sc. Ing. Bernardo Quiroga Turdera \quad \textbf{Duración:} 120 Minutos
\end{center}

\vspace{0.3cm}

\section*{1. Metadatos y Alineamiento Constructivo}
\begin{itemize}
    \item \textbf{Módulos:} 3 (Cierre) y 4 (Algoritmo Híbrido).
    \item \textbf{Resultados de Aprendizaje (ABET 1 y 6):} Formular matrices con impedancias complejas y modelar dinámicas frecuenciales generando visualizaciones logarítmicas de ingeniería.
    \item \textbf{Estrategia Didáctica:} Aprendizaje Basado en Problemas (ABP) enfocado en la transición de la estática (fasor único) a la dinámica espectral (barrido).
    \item \textbf{Entregables Activos:} Notebook 03 (v2), Diapositivas Beamer S3, Guía Analítica S3, Notebook 04.
\end{itemize}

\section*{2. Cronograma de Ejecución (120 Minutos)}
\begin{table}[h!]
    \centering
    \renewcommand{\arraystretch}{1.3}
    \begin{tabularx}{\textwidth}{@{} c c X l @{}}
        \toprule
        \textbf{Hora} & \textbf{Fase CDIO} & \textbf{Actividad Principal} & \textbf{Soporte} \\
        \midrule
        17:00 - 17:35 & \textbf{Implementar} & \textbf{El Gemelo Fasorial:} Matrices complejas (\texttt{1j}), \texttt{np.abs} y \texttt{np.linalg.solve} estático. & Notebook 03 \\
        17:35 - 17:55 & \textbf{Concebir} & \textbf{Límite Estático:} Reactancias dinámicas ($\omega$) y necesidad de escalas logarítmicas. & Beamer S3 \\
        17:55 - 18:15 & \textbf{Diseñar} & \textbf{El Límite Humano:} Cálculo a pulso de magnitudes en dos extremos frecuenciales. & Guía S3 \\
        18:15 - 18:45 & \textbf{Operar} & \textbf{Automatización:} Funciones modulares (\texttt{def}), \texttt{np.logspace} y ploteo (\texttt{plt.semilogx}). & Notebook 04 \\
        18:45 - 19:00 & \textbf{Evaluar} & \textbf{Diagnóstico:} Análisis visual de resonancia y cierre de la sesión. & Pizarra \\
        \bottomrule
    \end{tabularx}
\end{table}

\section*{3. Desarrollo Detallado por Fases}

\subsection*{Fase 1: Implementación Fasorial Estática (17:00 - 17:35)}
\textit{Objetivo: Demostrar que el álgebra lineal es una arquitectura universal (CC y CA).}
\begin{itemize}
    \item \textbf{Despliegue Notebook 03:} Los estudiantes declaran impedancias usando la constante nativa \texttt{1j}.
    \item \textbf{El Momento Eureka:} Resolver el sistema $\mathbf{Z}\mathbf{I} = \mathbf{V}$ usando exactamente la misma función \texttt{np.linalg.solve} de la sesión anterior, demostrando que Python maneja la trigonometría subyacente automáticamente.
\end{itemize}

\subsection*{Fase 2: El Quiebre Epistemológico (17:35 - 17:55)}
\textit{Objetivo: Justificar la entrada de la variable $\omega$ (Frecuencia) y la escala de Bode.}
\begin{itemize}
    \item \textbf{Intervención del Docente:} \textit{"El código que acaban de correr asume que el motor gira a velocidad constante. En la realidad, las señales aceleran y los sensores captan ruido. La impedancia $\mathbf{Z}$ mutará"}.
    \item \textbf{Clase Magistral:} Uso de las diapositivas Beamer S3 para introducir las funciones $X_L = \omega L$ y $X_C = -1/\omega C$. Explicar el concepto de "Década" logarítmica.
\end{itemize}

\subsection*{Fase 3: Misión de Ingeniería y Frustración Constructiva (17:55 - 18:15)}
\textit{Objetivo: Generar la necesidad inminente de la programación modular.}
\begin{itemize}
    \item \textbf{Cálculo Manual:} Resolver la Guía Analítica S3 calculando magnitudes para $\omega = 100$ y $\omega = 10,000$.
    \item \textbf{El Límite Humano:} Cuando la actividad consuma demasiado tiempo, pausar la clase y afirmar que para un análisis espectral se necesitan 1,000 puntos, lo que legitima el paso al Notebook 04.
\end{itemize}

\subsection*{Fase 4: Automatización y Dashboarding (18:15 - 19:00)}
\textit{Objetivo: Programar como ingenieros (Modularidad y Visualización).}
\begin{itemize}
    \item \textbf{Notebook 04:} Empaquetar la física en la función \texttt{def resolver\_red(w):}. Ejecutar un bucle \texttt{for} iterando sobre un vector de frecuencias \texttt{np.logspace}.
    \item \textbf{Detección de Resonancia:} Renderizar el gráfico con \texttt{plt.semilogx}. Realizar la evaluación socrática final: \textit{"En el pico de la gráfica, las reactancias del capacitor y la bobina se anularon mutuamente. ¿Qué riesgos implica esto en el encendido de un actuador mecatrónico real?"}
\end{itemize}

\end{document}
